% Generated draft: requires mathematical and source review.
\documentclass[11pt]{article}
\usepackage[margin=25mm]{geometry}
\usepackage{fontspec}
\setmainfont{Latin Modern Roman}
\usepackage{amsmath,amssymb,mathtools,unicode-math}
\setmathfont{Latin Modern Math}
\usepackage{enumitem,needspace,xurl,hyperref}
\setlength{\parindent}{0pt}
\setlength{\parskip}{5pt}
\setlength{\emergencystretch}{3em}
\newcommand{\N}{\mathbb{N}}
\newcommand{\Z}{\mathbb{Z}}
\newcommand{\Q}{\mathbb{Q}}
\newcommand{\R}{\mathbb{R}}
\newcommand{\C}{\mathbb{C}}
\newcommand{\F}{\mathbb{F}}
\newcommand{\A}{\mathbb{A}}
\newcommand{\PP}{\mathbb{P}}
\newcommand{\E}{\mathbb{E}}
\newcommand{\eps}{\varepsilon}
\newcommand{\dd}{\,\mathrm{d}}
\begin{document}
\section*{UnsolvedMath ID: 10900045}
\textbf{Source problem code:} catalogue entry
\subsection{Definitions and mathematical statement}
Let $S$ be a surface in the intended mapping-class-group setting, and let $\operatorname{Mod}(S)$ denote isotopy classes of its self-homeomorphisms, with orientation and marking conventions to be specified. Given a probability measure $\nu$ on this group, a random walk started at the identity is
\[
 W_n=X_1\cdots X_n,\qquad X_1,X_2,\ldots\ \text{independent with law }\nu.
\]
A pseudo-Anosov mapping class has transverse invariant stable and unstable measured foliations, scaled by reciprocal factors $\lambda$ and $\lambda^{-1}$ with $\lambda>1$. An interior foliation singularity is \emph{trivalent} when it has exactly three prongs. For a closed orientable surface, the generic principal-stratum interpretation asks that all singularities of both invariant foliations have three prongs. Punctures and boundary components require separate conventions.

\textbf{Problem (I. Kapovitch).} Show that the random walk gives a pseudo-Anosov mapping class with generic trivalent invariant-foliation singularities with probability tending to one as its length tends to infinity. With the closed-surface interpretation just described, the requested limit is
\[
 \lim_{n\to\infty}\Pr_\nu\bigl(W_n\text{ is pseudo-Anosov and every singularity of its stable and unstable foliations is three-pronged}\bigr)=1.
\]
The extracted statement does not specify $S$, the law $\nu$, its support, moment or generation hypotheses, or whether the assertion is for a particular simple random walk or for a class of walks. Thus the display is a precise target once those data are supplied, not a claim for every probability law on every mapping class group. In particular, the needed nondegeneracy assumptions cannot be omitted. If punctures are allowed, the source must say whether ``trivalent'' concerns only interior singularities and what prongs are permitted at punctures.

The pseudo-Anosov conclusion alone is weaker than the asserted singularity-pattern conclusion. The registry summary attributes probability tending to one for the former to random-walk results and leaves the latter unverified.
\subsection{Short English statement}
Prove that a sufficiently long mapping-class-group random walk produces a pseudo-Anosov whose invariant foliations have only the intended generic three-pronged singularities, with probability tending to one. The surface and random-walk hypotheses are missing from the excerpt.
\subsection{Sources}
\begin{itemize}
\item Local UnsolvedMath registry, permanent ID 10900045, source code catalogue entry; attributed exactly as supplied to I. Kapovitch, no proposal year supplied.
\item Delp, Hoffoss and Manning, \emph{Problems In Groups, Geometry, and Three-Manifolds} (2015), subsection 6.1 (I. Kapovitch), PDF page 5; \url{https://arxiv.org/pdf/1512.04620#page=5} and \url{https://arxiv.org/abs/1512.04620}. The provenance records source-TeX extraction checked against rendered PDF text and access on 2026-07-29.
\item The research summary refers to Maher, ``Exponential decay in the mapping class group,'' and related observations of Rivin on random Heegaard/gluing maps. It supplies no precise bibliographic identifiers or theorem hypotheses for these references.
\end{itemize}
\end{document}
